\chapter{Qualitative Study}

\section{Introduction}

Consider the discrete dynamical system
\begin{equation} \label{DDS}
x_{i+1} = f(x_i) \ .
\end{equation}
For a qualitative study of this system, we would like to find all the
fixed points, periodic orbits, \ldots.  We would also like to find the
sets of points forward asymptotic to these objects.

We first study the fixed points of (\ref{DDS}).

\begin{defn}
Let $p$ be a fixed point of $f$.  $p$ is {\bfseries
hyperbolic}\index{Fixed point!Hyperbolic fixed point} if
$\left|f'(p)\right| \neq 1$.
\end{defn}

A proof similar to the proof of the Fixed Point Theorem shows that:

\begin{prop} \label{fp_stability}
Let $p$ be an hyperbolic fixed point of $f$.  If
$\left|f'(p)\right|<1$, then there exists an open interval $I$
containing $p$ and such that
\[
\lim_{j\rightarrow \infty} f^j(x) = p
\]
for all $x \in I$.
However, if $\left|f'(p)\right|>1$, then there exists an open
interval $I$ containing $p$ such that for each $x$ in $I$,
$x \neq p$, one can find $j \in \NN$ such that $f^j(x) \not\in I$.
\end{prop}

\section{Logistic Equation}\index{Logistic equation}

This section is about the logistic equation
\[
x_{i+1} = f_\mu(x_i) = \mu x_i (1-x_i) \ .
\]
However, most of the observations about the logistic map are true for
similarly shaped mappings of $I=[0,1]$ into $I$.

For $\mu<3$, all iterations converge to the sink at $p_\mu$.   For
$\mu$ slightly bigger than $3$, the fixed point $p_\mu$ becomes a
source and there appears an attracting  periodic orbit of period $2$
--- All orbits eventually ``bounce back and for'' between the two
points of the orbit.    We say that there is a
{\bfseries bifurcation point}\index{Bifurcation point} at $\mu=3$.

When $\mu \approx 3.236\ldots$, the periodic orbit of period $2$
becomes unstable and there appears an attracting periodic orbit of
period $4$.  Then, there is a bifurcation from a attracting periodic
orbit of period $4$ to an attracting periodic orbit of period $8$,
And so on.  This is the best known example of a
{\bfseries period doubling cascade}\index{Period doubling}.

Period doubling accumulates to $\mu = 3.5699456\ldots$.  This number
is known as the {\em Feigenbaum point}.  Moreover, let $\mu_0=3$,
$\mu_1=3.236\ldots$, $\mu_2$, $\mu_3$, \ldots be the values of $\mu$
for which the logistic mapping undergoes period doubling and let
$d_j = \mu_{j+1}-\mu_j$.  It has been showed that
\[
\lim_{j\rightarrow \infty} \frac{d_j}{d_{j+1}} =
4.6692016091029\ldots
\]
This number is called the
{\em Feigenbaum constant}\index{Feigenbawm number}.  Feigenbaum
discovered this number in 1975.  This constance is universal in the
sense that it is the same for a whole class of iterative systems of
the form $x_{n+1} = g_\mu(x_n)$ where the graph of $g_\mu(x)$ looks
like the graph of $f_\mu(x)$.

Period doubling is far from being the must complex type of
bifurcation. To understand the complex behaviour of the orbits of
$f_\mu$, we need the following theorem.

\begin{theo}[Sarkovskii]\index{Sarkovskii's Theorem}
Consider the order on the positive integers
\nomenclature{$\gg$}{Sarkovskii's ordering} defined by
\begin{multline*}
3 \gg 5 \gg 7 \gg \ldots
\gg 3\times 2 \gg 5\times 2 \gg \ldots \\
\gg 3\times 2^2 \gg 5\times 2^2 \gg \ldots
\gg 3\times 2^3 \gg 5\times 2^3 \gg \ldots
2^4 \gg 2^3 \gg 2^2 \gg 2 \gg 1 \quad .
\end{multline*}
Let $f:\RR \rightarrow \RR$ be a continuous function and $k$ be a
(prime) number.  If $f$ has a periodic point of period $k$, then $f$ has
a periodic point of period $m$ for all $m \ll k$.
\end{theo}

\begin{proof}
A proof of this theorem is given in \cite[page 60]{D}.
\end{proof}

This theorem is often stated as period three leads to chaos.
A definition of chaos is given in \cite[Definition 8.5]{D}
The definition given in \cite{D} is not the most economic, a simpler
definition is given in \cite{BBCDS, AG} 

%%% Local Variables: 
%%% mode: latex
%%% TeX-master: "template"
%%% End: 
